\documentclass[12pt,a4paper]{article}
\usepackage[utf8]{inputenc}
\usepackage[T5]{fontenc}
\usepackage[vietnamese]{babel}
\usepackage{amsmath, amssymb, amsfonts}
\usepackage{geometry}
\geometry{margin=2.5cm}
\newcommand{\R}{\mathbb{R}}
\newcommand{\N}{\mathbb{N}}

\title{Buổi 07 -- Ngày 06-09-2026 -- môn Đại số tuyến tính -- lớp MA003.F32.LT.CNTT}
\date{}
\author{}

\begin{document}
\maketitle

\section*{8/ MA TRẬN CHUYỂN CƠ SỞ}

Cho $(V,+,\cdot)$ là không gian tuyến tính tổng quát trên $F=\R$

và cho $a=\{a_1,a_2,\ldots,a_n\}$ là một cơ sở của $V$;

$b=\{b_1,b_2,\ldots,b_n\}$ là một cơ sở của $V$.

Ta lấy tọa độ của các véc tơ trong $b$ biểu diễn theo cơ sở $a$, cụ thể như sau:

$[b_1]_a$; $[b_2]_a$; \ldots; $[b_n]_a$.

(ta phải giải $n$ hệ pt, mỗi hệ gồm $n$ pt và $n$ ẩn số)

Ta lập ma trận $P=T=T_{a\to b}=T_{ab}=\big([b_1]_a \ [b_2]_a \ \cdots \ [b_n]_a\big)$ ma trận vuông cấp $n$.

= ma trận chuyển cơ sở từ $a$ sang $b$.

\medskip
\noindent\textbf{\textit{* Tính chất}}:

Cho $a=\{a_1,a_2,\ldots,a_n\}$; $b=\{b_1,b_2,\ldots,b_n\}$; $\varsigma=c=\{c_1,c_2,\ldots,c_n\}$ là các cơ sở của không gian tuyến tính $V$. Ta có một số tính chất sau:

i/ $T_{a\to a}=T_{b\to b}=T_{\varsigma\to\varsigma}=I_n$.

ii/ $T_{b\to a}=T_{ba}=(T_{ab})^{-1}=(T_{a\to b})^{-1}$.

iii/ $T_{a\to b}=T_{a\to c}\cdot T_{c\to b}$

\hfill khó tìm \hfill dễ tìm

Ta thường chọn $c=\beta_0$ cơ sở chính tắc của không gian $V$; khi đó ta viết

\[
T_{a\to b}=T_{a\to \beta_0}\cdot T_{\beta_0\to b}=(T_{\beta_0\to a})^{-1}\cdot T_{\beta_0\to b}
\]

\section*{9/ CÔNG THỨC ĐỔI TỌA ĐỘ}

Ta có
\[
[\alpha]_a=T_{a\to b}\cdot[\alpha]_b=T_{a\to c}\cdot[\alpha]_c=T_{a\to\beta_0}\cdot[\alpha]_{\beta_0}=T_{a\to d}\cdot[\alpha]_d
\]

\noindent\rule{\linewidth}{0.4pt}

\noindent\underline{Ví dụ mẫu 6}: bài 3.31

\medskip
\noindent\textbf{Bài 3.31.} Trong không gian tuyến tính $\R^3$ cho hai hệ cơ sở $(a)=\{a_1,a_2,a_3\}$ và $(b)=\{b_1,b_2,b_3\}$ với
\[
a_1=(2,-1,3),\ a_2=(1,1,2),\ a_3=(2,1,4),
\]
\[
b_1=(1,2,3),\ b_2=(3,1,-2),\ b_3=(-1,1,2).
\]
Hãy tính ma trận chuyển cơ sở từ hệ $(a)$ sang hệ $(b)$.

\medskip
\noindent\underline{Giải}:

\noindent\underline{Cách 1}: Dùng định nghĩa, ta có:
\[
T_{a\to b}=T_{ab}=\big([b_1]_a \ \ [b_2]_a \ \ [b_3]_a\big)
\]

Ta tìm $[b_1]_a$ như sau

Ta tìm $c_1,c_2,c_3\in\R$ sao cho $b_1=c_1a_1+c_2a_2+c_3a_3$

\[
\Leftrightarrow (1,2,3)=c_1(2,-1,3)+c_2(1,1,2)+c_3(2,1,4)
\]
\[
\Leftrightarrow
\begin{cases}
2c_1+c_2+2c_3=1\\
-c_1+c_2+c_3=2\\
3c_1+2c_2+4c_3=3
\end{cases}
\Leftrightarrow
\begin{cases}
c_1=-1\\
c_2=-1\\
c_3=2
\end{cases}
\]

Ta tìm $[b_2]_a$ như sau

Ta tìm $c_1,c_2,c_3\in\R$ sao cho $b_2=c_1a_1+c_2a_2+c_3a_3$

\[
\Leftrightarrow (3,1,-2)=c_1(2,-1,3)+c_2(1,1,2)+c_3(2,1,4)
\]
\[
\Leftrightarrow
\begin{cases}
2c_1+c_2+2c_3=3\\
-c_1+c_2+c_3=1\\
3c_1+2c_2+4c_3=-2
\end{cases}
\Leftrightarrow
\begin{cases}
c_1=8\\
c_2=31\\
c_3=-22
\end{cases}
\]

Ta tìm $[b_3]_a$ như sau

Ta tìm $c_1,c_2,c_3\in\R$ sao cho $b_3=c_1a_1+c_2a_2+c_3a_3$

\[
\Leftrightarrow (-1,1,2)=c_1(2,-1,3)+c_2(1,1,2)+c_3(2,1,4)
\]
\[
\Leftrightarrow
\begin{cases}
2c_1+c_2+2c_3=-1\\
-c_1+c_2+c_3=1\\
3c_1+2c_2+4c_3=2
\end{cases}
\Leftrightarrow
\begin{cases}
c_1=-4\\
c_2=-13\\
c_3=10
\end{cases}
\]

Ta có ma trận chuyển cơ sở cần tìm là:
\[
T_{ab}=\big([b_1]_a \ \ [b_2]_a \ \ [b_3]_a\big)=
\begin{pmatrix}
-1 & 8 & -4\\
-1 & 31 & -13\\
2 & -22 & 10
\end{pmatrix}
\]

\noindent\underline{Cách 2}: Ta có: $T_{ab}=T_{a\to b}=T_{a\to\beta_0}\cdot T_{\beta_0\to b}$

Ta có: $T_{\beta_0\to a}=\big([a_1]_{\beta_0}\ \ [a_2]_{\beta_0}\ \ [a_3]_{\beta_0}\big)=
\begin{pmatrix}
2 & 1 & 2\\
-1 & 1 & 1\\
3 & 2 & 4
\end{pmatrix}$

và $T_{\beta_0\to b}=\big([b_1]_{\beta_0}\ \ [b_2]_{\beta_0}\ \ [b_3]_{\beta_0}\big)=
\begin{pmatrix}
1 & 3 & -1\\
2 & 1 & 1\\
3 & -2 & 2
\end{pmatrix}$

\begin{align*}
\Rightarrow T_{ab}=T_{a\to b} &=
\begin{pmatrix}
2 & 1 & 2\\
-1 & 1 & 1\\
3 & 2 & 4
\end{pmatrix}^{-1}
\cdot
\begin{pmatrix}
1 & 3 & -1\\
2 & 1 & 1\\
3 & -2 & 2
\end{pmatrix}
\\
&=
\begin{pmatrix}
2 & 0 & -1\\
7 & 2 & -4\\
-5 & -1 & 3
\end{pmatrix}
\cdot
\begin{pmatrix}
1 & 3 & -1\\
2 & 1 & 1\\
3 & -2 & 2
\end{pmatrix}
=
\begin{pmatrix}
-1 & 8 & -4\\
-1 & 31 & -13\\
2 & -22 & 10
\end{pmatrix}
\end{align*}

\medskip
\noindent\underline{Ví dụ mẫu 7}: bài 3.29

\medskip
\noindent\textbf{Bài 3.29.} Trong không gian tuyến tính ba chiều $U$ cho hai hệ cơ sở $(a)$ và $(b)$ với ma trận chuyển cơ sở từ hệ $(a)$ sang hệ $(b)$ là
\[
T=\begin{pmatrix}4 & 2 & 1\\ 1 & -2 & 3\\ 3 & 3 & 4\end{pmatrix}.
\]
Hãy tính ma trận chuyển cơ sở từ hệ $(b)$ sang hệ $(a)$.

\medskip
\noindent\underline{Giải}:

Ta có ma trận chuyển cơ sở từ $a$ sang $b$ là:
\[
T_{ab}=\begin{pmatrix}4 & 2 & 1\\ 1 & -2 & 3\\ 3 & 3 & 4\end{pmatrix}
\]

Nên ma trận chuyển cơ sở từ $b$ sang $a$ là:
\[
T_{ba}=(T_{ab})^{-1}=\begin{pmatrix}4 & 2 & 1\\ 1 & -2 & 3\\ 3 & 3 & 4\end{pmatrix}^{-1}
=\begin{pmatrix}17/49 & 5/49 & -8/49\\ -5/49 & -13/49 & 11/49\\ -9/49 & 6/49 & 10/49\end{pmatrix}
=\frac{1}{49}\begin{pmatrix}17 & 5 & -8\\ -5 & -13 & 11\\ -9 & 6 & 10\end{pmatrix}
\]

\medskip
\noindent\textit{\underline{Bài tập tương tự}}: bài 3.28; 3.30; 3.32; 3.33

\medskip
\noindent\textbf{Bài 3.28.} Trong không gian tuyến tính ba chiều $U$ cho hai hệ cơ sở $(a)$ và $(b)$ với ma trận chuyển cơ sở từ hệ $(a)$ sang hệ $(b)$ là
\[
T=\begin{pmatrix}1 & 2 & -4\\ 2 & 5 & 3\\ 3 & 2 & 1\end{pmatrix}.
\]
Cho biết phần tử $x$ có tọa độ trong cơ sở thứ nhất $(a)$ là $x_a=(1,4,-2)$. Hãy tính tọa độ $x_b$ của phần tử $x$ trong cơ sở thứ hai $(b)$.

\medskip
\noindent\textbf{Bài 3.30.} Trong không gian tuyến tính ba chiều $U$ cho hai hệ cơ sở $(a)=\{a_1,a_2,a_3\}$ và $(b)=\{b_1,b_2,b_3\}$ với
\[
b_1=2a_1+3a_2-a_3,\quad b_2=a_1+4a_2+2a_3,\quad b_3=3a_1-a_2+a_3.
\]
Hãy tính ma trận chuyển cơ sở từ hệ $(b)$ sang hệ $(a)$.

\medskip
\noindent\textbf{Bài 3.32.} Trong không gian tuyến tính ba chiều $U$ cho ba hệ cơ sở $(e)$, $(a)$ và $(b)$. Cho biết ma trận chuyển cơ sở từ cơ sở $(e)$ sang cơ sở $(a)$ là
\[
T_{ea}=\begin{pmatrix}2 & -1 & 1\\ 2 & 1 & 2\\ 3 & 1 & 4\end{pmatrix}
\]
và ma trận chuyển cơ sở từ cơ sở $(e)$ sang cơ sở $(b)$ là
\[
T_{eb}=\begin{pmatrix}1 & 1 & 1\\ -2 & 3 & 1\\ 2 & 1 & -2\end{pmatrix}.
\]
Hãy tính ma trận chuyển cơ sở từ cơ sở $(a)$ sang cơ sở $(b)$.

\medskip
\noindent\textbf{Bài 3.33.} Trong không gian tuyến tính $\R^3$ cho hai hệ cơ sở $(a)=\{a_1,a_2,a_3\}$ và $(b)=\{b_1,b_2,b_3\}$ với
\[
a_1=(3,1,4),\ a_2=(5,-4,2),\ a_3=(2,1,1),
\]
\[
b_1=(3,-2,3),\ b_2=(4,1,-2),\ b_3=(3,4,2).
\]
Hãy tính ma trận chuyển cơ sở từ hệ $(b)$ sang hệ $(a)$.

\newpage
\section*{CHƯƠNG 4: TRỊ RIÊNG, VÉC TƠ RIÊNG,\\ KHÔNG GIAN RIÊNG, ĐA THỨC ĐẶC TRƯNG CỦA MA TRẬN VUÔNG}

\section*{1/ MỘT SỐ KHÁI NIỆM}

Cho $A$ là ma trận vuông cấp $n$ trên trường số thực $\R$.

Cho số thực $c\in\R$.

Đặt $E_c=\{X\in\R^n \mid (A-cI_n)X=O\}$ không gian nghiệm của hệ pt tuyến tính thuần nhất $(A-cI_n)X=O$.

Ta giải hệ pt tuyến tính thuần nhất $(A-cI_n)X=O$.

Nếu hệ pt này chỉ có 1 bộ nghiệm tầm thường duy nhất là $(0,0,0,\ldots,0)\Rightarrow E_c=\{(0;0;0;\ldots;0)\}$ thì $c$ không là một trị riêng của $A$.

Nếu hệ pt này có vô số nghiệm (bên cạnh bộ nghiệm tầm thường $(0,0,0,\ldots,0)$) thì ta nói $c$ là một trị riêng (eigenvalue) của $A$, và viết $E_c\neq\{(0;0;0;\ldots;0)\}$.

Khi đó, ta gọi $E_c=$ không gian riêng (eigenvector-space), ứng với trị riêng $c$.

Và mỗi véc tơ $\alpha\in E_c$ ta gọi là véc tơ riêng (eigenvector), ứng với trị riêng $c$.

\medskip
\noindent\underline{Ví dụ mẫu 1}: Cho ma trận thực

\[
A=\begin{pmatrix}-3 & -1\\ 7 & 5\end{pmatrix}
\]

a/ Hỏi $c=3$ có phải là 1 trị riêng của $A$ không? Vì sao?

b/ Hỏi $c=-2$ có phải là 1 trị riêng của $A$ không? Vì sao?

c/ Hỏi $c=4$ có phải là 1 trị riêng của $A$ không? Vì sao?

\medskip
\noindent\underline{Giải}:

Ta có $c$

a/ Với $c=3$ ta có $E_c=\{X\in\R^2 \mid (A-cI_2)X=O\}$

Giải hệ pt $E_3=\{X\in\R^2 \mid (A-3I_2)X=O\}$, tức $(A-3I_2)X=O$

\[
\left[\begin{pmatrix}-3 & -1\\ 7 & 5\end{pmatrix}-3\begin{pmatrix}1 & 0\\ 0 & 1\end{pmatrix}\ \middle|\ \begin{matrix}0\\0\end{matrix}\right]
\Leftrightarrow
\left[\begin{pmatrix}-6 & -1\\ 7 & 2\end{pmatrix}\ \middle|\ \begin{matrix}0\\0\end{matrix}\right]
\xrightarrow{(2)\to 6(2)+7(1)}
\begin{bmatrix}-6 & -1\\ 0 & 5\end{bmatrix}\ \big|\ \begin{matrix}0\\0\end{matrix}
\]

Ta có:
\[
\begin{cases}-6x_1-x_2=0\\ 5x_2=0\end{cases}
\Rightarrow
\begin{cases}x_1=0\\ x_2=0\end{cases}
\]

\[
\Rightarrow E_c=E_3=\{X=(x_1,x_2)=(0,0)\}
\]

là bộ nghiệm tầm thường nên $c=3$ không phải là trị riêng của $A$.

b/ Với $c=-2$ ta có $E_{-2}=\{X\in\R^2 \mid (A-(-2)I_2)X=O\}$

Giải hệ pt $(A+2I_2)X=O$

\[
\left[\begin{pmatrix}-3 & -1\\ 7 & 5\end{pmatrix}+2\begin{pmatrix}1 & 0\\ 0 & 1\end{pmatrix}\ \middle|\ \begin{matrix}0\\0\end{matrix}\right]
\Leftrightarrow
\left[\begin{pmatrix}-1 & -1\\ 7 & 7\end{pmatrix}\ \middle|\ \begin{matrix}0\\0\end{matrix}\right]
\xrightarrow{(2)\to(2)+7(1)}
\begin{bmatrix}-1 & -1\\ 0 & 0\end{bmatrix}\ \big|\ \begin{matrix}0\\0\end{matrix}
\]

Do cột 2 không bán chuẩn hóa được nên hệ có vô số nghiệm như sau:

Đặt $x_2=t,\ \forall t\in\R$.

Ta có:
\[
-x_1-x_2=0\Rightarrow x_1=-x_2=-t
\]
\[
\Rightarrow E_c=E_{-2}=\{X=(x_1,x_2)=(-t,t)\mid t\in\R\}.
\]

là bộ nghiệm tầm thường nên $c=-2$ là trị riêng của $A$.

c/ Với $c=4$ ta có $E_4=\{X\in\R^2 \mid (A-4I_2)X=O\}$, tức $(A-4I_2)X=O$

\medskip
\noindent\textit{\underline{Bài tập tương tự}}:

a/ $A=\begin{pmatrix}4 & -2\\ -11 & -5\end{pmatrix}$. Hỏi $c=-5$; $c=-7$ có phải là trị riêng của $A$ không?

b/ $A=\begin{pmatrix}8 & 2\\ -4 & 2\end{pmatrix}$. Hỏi $c=0$; $c=6$ có phải là trị riêng của $A$ không?

c/ $A=\begin{pmatrix}5 & 6 & -3\\ -1 & 0 & 1\\ 1 & 2 & 1\end{pmatrix}$. Hỏi $c=-3$; $c=2$ có phải là trị riêng của $A$ không?

d/ $A=\begin{pmatrix}7 & -6 & -10\\ -12 & 17 & 24\\ 12 & -15 & -22\end{pmatrix}$. Hỏi $c=0$; $c=-1$ có phải là trị riêng của $A$ không?

\newpage
\section*{2/ ĐA THỨC ĐẶC TRƯNG}

Cho $A$ là ma trận vuông cấp $n$ trên $\R$. Đa thức đặc trưng của ma trận $A$ được định nghĩa:
\[
p_A(x)=\det(xI_n-A)
\]
hoặc
\[
p_A(\lambda)=\det(A-\lambda I_n)
\]

Số $c$ là trị riêng của $A$ khi và chỉ khi $p_A(c)=0$, tức là $c$ là nghiệm của phương trình đặc trưng $p_A(x)=0$ (hoặc $p_A(\lambda)=0$).

\medskip
\noindent\underline{Ví dụ mẫu}: Cho ma trận thực $A=\begin{pmatrix}-3 & -1\\ 7 & 5\end{pmatrix}$.

Tìm tất cả các trị riêng của $A$.

\medskip
\noindent\underline{Giải}:

\[
p_A(x)=\det(xI_2-A)=\det\left[x\begin{pmatrix}1&0\\0&1\end{pmatrix}-\begin{pmatrix}-3&-1\\7&5\end{pmatrix}\right]
=\det\begin{pmatrix}x+3 & 1\\ -7 & x-5\end{pmatrix}
\]
\[
=(x+3)(x-5)-1\cdot(-7)=x^2-2x-15+7=x^2-2x-8
\]
\[
p_A(x)=0\Leftrightarrow x^2-2x-8=0\Leftrightarrow (x-4)(x+2)=0\Leftrightarrow
\begin{bmatrix}x=4\\ x=-2\end{bmatrix}
\]
\[
\Rightarrow \begin{cases}c_1=\lambda_1=4\\ c_2=\lambda_2=-2\end{cases}
\]

Vậy ma trận $A$ có 2 trị riêng là $4$ và $-2$.

\newpage
\section*{3/ SỐ MŨ CỦA CÁC TRỊ RIÊNG}

Cho ma trận vuông $A$ cấp $n$ trên $\R$. Giả sử đa thức đặc trưng $p_A(\lambda)=0$ có nghiệm được phân tích thành:
\[
x_1=\lambda_1=\ldots,\quad x_2=\lambda_2=\ldots,\quad x_3=\lambda_3=\ldots
\]

hoặc tổng quát, đa thức đặc trưng $p_A(x)$ (hay $p_A(\lambda)$) được phân tích thành nhân tử:
\[
p_A(x)=(x-c_1)^{r_1}(x-c_2)^{r_2}\cdots(x-c_k)^{r_k}
\]
\[
\begin{cases}
c_1=\lambda_1=\ldots;\ r_1=\ldots\\
c_2=\lambda_2=\ldots;\ r_2=\ldots\\
\ldots\ldots\ldots\ldots\ldots\ldots\ldots\ldots\ldots\\
c_k=\lambda_k=\ldots;\ r_k=\ldots
\end{cases}
\]

Ví dụ: $p_A(\lambda)=(\lambda-4)^2(\lambda+5)\Rightarrow
\begin{cases}\lambda_1=4;\ r_1=2\\ \lambda_2=-5;\ r_2=1\end{cases}$

\medskip
\noindent(Trường hợp đặc biệt $r_1=r_2=\cdots=r_k=1$: mỗi trị riêng có số mũ bằng 1.)

\newpage
\section*{4/ CHIỀU CỦA KHÔNG GIAN RIÊNG}

Cho ma trận vuông $A$ cấp $n$ trên $\R$; gọi $a_j$ là cơ sở của không gian riêng
\[
E_{c_j}=E_{\lambda_j}=\{X\in\R^n \mid (A-c_jI_n)X=O\},\quad 1\le j\le k,\ j\in\{1,2,\ldots,k\}.
\]

Ta có định lý: $\dim_\R E_{c_j}=\cdots$? Câu trả lời:
\[
\dim_\R E_{c_j}<r_j \quad\text{hoặc}\quad \dim_\R E_{c_j}=r_j,
\]
với chú ý đặc biệt khi $r_j>1$.

Đặt $a=a_1\cup a_2\cup\cdots\cup a_k$.

\newpage
\section*{5/ CHÉO HÓA MA TRẬN VUÔNG}

Nếu $a=a_1\cup a_2\cup\cdots\cup a_k$ là một cơ sở của $\R^n$ (tức là $\dim_\R E_{c_j}=r_j$ với mọi $j$), gọi $\beta_0$ là cơ sở chính tắc, đặt
\[
P=T=T_{\beta_0\to a}=\cdots
\]

Lúc này $T$ luôn khả nghịch, nghĩa là luôn có $T^{-1}$ sao cho:
\[
D=T^{-1}AT=
\begin{pmatrix}
c_1 & & & & &\\
 & \ddots & & & &\\
 & & c_1 & & &\\
 & & & c_2 & &\\
 & & & & \ddots &\\
 & & & & & c_k
\end{pmatrix}
\]

+ Ta gọi đây là dạng chéo hóa cần tìm của $A$.

\medskip
\noindent\textbf{\textit{* Ứng dụng}}: tìm lũy thừa $A^m$, với mọi $m\ge 0,\ m\in\N$.

Ta có: $T^{-1}AT=\operatorname{diag}(c_1,\ldots,c_1,c_2,\ldots,c_2,\ldots,c_k,\ldots,c_k)\Rightarrow A=T\cdot\operatorname{diag}(c_1,\ldots,c_k)\cdot T^{-1}$.

Nhân vế theo vế $m$ lần ta có:
\[
A\cdot A\cdot A\cdots A=T\operatorname{diag}(c_1,\ldots,c_k)T^{-1}\cdot T\operatorname{diag}(c_1,\ldots,c_k)T^{-1}\cdots T\operatorname{diag}(c_1,\ldots,c_k)T^{-1}
\]
\[
\Rightarrow A^m=T\operatorname{diag}(c_1,\ldots,c_k)^m T^{-1}=T\operatorname{diag}((c_1)^m,\ldots,(c_k)^m)T^{-1},\quad \forall m\in\N.
\]

(Ta phải tìm $T^{-1}$ và thực hiện phép nhân ma trận để tìm $A^m$).

\medskip
\noindent\underline{Ví dụ mẫu}: Cho ma trận thực $A=\begin{pmatrix}3 & -1 & 0\\ -1 & 2 & -1\\ 0 & -1 & 3\end{pmatrix}$.

Hãy chéo hóa $A$, sau đó tìm $A^m$, $\forall m\in\N$.

\medskip
\noindent\underline{Giải}:

Ta có đa thức đặc trưng của $A$:
\[
p_A(\lambda)=\det(A-\lambda I_3)=\det\left(
\begin{pmatrix}3 & -1 & 0\\ -1 & 2 & -1\\ 0 & -1 & 3\end{pmatrix}
-\lambda\begin{pmatrix}1 & 0 & 0\\ 0 & 1 & 0\\ 0 & 0 & 1\end{pmatrix}\right)
\]
\[
=\begin{vmatrix}3-\lambda & -1 & 0\\ -1 & 2-\lambda & -1\\ 0 & -1 & 3-\lambda\end{vmatrix}
\]
\[
=(3-\lambda)^2(2-\lambda)+(-1)(-1)\cdot 0+0\cdot(-1)(-1)-\big[0\cdot(2-\lambda)+(-1)(-1)(3-\lambda)+(3-\lambda)(-1)(-1)\big]
\]
\[
=(3-\lambda)^2(2-\lambda)-(3-\lambda)-(3-\lambda)
=(3-\lambda)\big[(3-\lambda)(2-\lambda)-2\big]
=(3-\lambda)\big[6-5\lambda+\lambda^2-2\big]
\]
\[
=(3-\lambda)(\lambda^2-5\lambda+4)=(3-\lambda)(\lambda-1)(\lambda-4)
\]

Giải pt:
\[
p_A(\lambda)=0\Leftrightarrow (3-\lambda)(\lambda-1)(\lambda-4)=0\Leftrightarrow
\begin{bmatrix}\lambda=3\\ \lambda=1\\ \lambda=4\end{bmatrix}
\]

Do đó, ta nói $A$ có 3 trị riêng: $\lambda_1=3;\ \lambda_2=1;\ \lambda_3=4$ và có số mũ của các trị riêng là:
\[
r_1=1;\ r_2=1;\ r_3=1
\]

Tiếp theo, ta tìm cơ sở $a_j$ cho các không gian riêng:

$E_{c_j}=E_{\lambda_j}=\{X\in\R^3 \mid (A-\lambda_j I_3)X=O\}$, với $1\le j\le 3$.

\medskip
Với $\lambda_1=3$, ta có:

$E_{\lambda_1}=E_3=\{X\in\R^3 \mid (A-\lambda_1 I_3)X=O\}$

Ta có dạng ma trận hóa của hệ pt: $(A-\lambda_1I_3)X=O$

\[
\Leftrightarrow
\left[\begin{pmatrix}3 & -1 & 0\\ -1 & 2 & -1\\ 0 & -1 & 3\end{pmatrix}-3\begin{pmatrix}1&0&0\\0&1&0\\0&0&1\end{pmatrix}\ \middle|\ \begin{matrix}0\\0\\0\end{matrix}\right]
\]
\[
\Leftrightarrow
\begin{bmatrix}0 & -1 & 0\\ -1 & -1 & -1\\ 0 & -1 & 0\end{bmatrix}\big|\begin{matrix}0\\0\\0\end{matrix}
\xrightarrow{(1)\leftrightarrow(2)}
\begin{bmatrix}-1 & -1 & -1\\ 0 & -1 & 0\\ 0 & -1 & 0\end{bmatrix}\big|\begin{matrix}0\\0\\0\end{matrix}
\xrightarrow[(3)\to(3)-(2)]{(1)\to -1(1)}
\begin{bmatrix}1 & 1 & 1\\ 0 & -1 & 0\\ 0 & 0 & 0\end{bmatrix}\big|\begin{matrix}0\\0\\0\end{matrix}
\]

Do cột 3 không bán chuẩn hóa được nên hệ có vô số nghiệm như sau:

Đặt $x_3=a,\ \forall a\in\R$.

Ta có:
\[
\begin{cases}x_1+x_2+x_3=0\\ -x_2=0\end{cases}
\Rightarrow
\begin{cases}x_1=-x_2-x_3=-a\\ x_2=0\end{cases}
\]

\[
\Rightarrow E_{\lambda_1}=E_3=\{X=(x_1,x_2,x_3)=(-a,0,a)\mid a\in\R\}=\{X=a(-1,0,1)\mid a\in\R\}
\]

Đặt $u_1=(-1,0,1)\neq(0,0,0)=O$

$u_1=(-1,0,1)$ là tập sinh của $E_3$.

Ta có: $u_1=(-1,0,1)$ là một véc tơ riêng; và $a_1=S_1=\{u_1\}$ là một cơ sở của $E_3$

và số chiều là $\dim_\R E_{\lambda_1}=1=r_1$ \hfill (1)

\medskip
Tương tự,

Với $\lambda_2=1$:

$E_{\lambda_2}=E_1=\{X\in\R^3 \mid (A-\lambda_2I_3)X=O\}$

Ta có dạng ma trận hóa của hệ pt tuyến tính: $(A-\lambda_2I_3)X=O$

\[
\Leftrightarrow
\left[\begin{pmatrix}3 & -1 & 0\\ -1 & 2 & -1\\ 0 & -1 & 3\end{pmatrix}-1\begin{pmatrix}1&0&0\\0&1&0\\0&0&1\end{pmatrix}\ \middle|\ \begin{matrix}0\\0\\0\end{matrix}\right]
\]
\[
\Leftrightarrow
\begin{bmatrix}2 & -1 & 0\\ -1 & 1 & -1\\ 0 & -1 & 2\end{bmatrix}\big|\begin{matrix}0\\0\\0\end{matrix}
\xrightarrow{(2)\to 2(2)+(1)}
\begin{bmatrix}2 & -1 & 0\\ 0 & 1 & -2\\ 0 & -1 & 2\end{bmatrix}\big|\begin{matrix}0\\0\\0\end{matrix}
\xrightarrow{(3)\to(3)+(2)}
\begin{bmatrix}2 & -1 & 0\\ 0 & 1 & -2\\ 0 & 0 & 0\end{bmatrix}\big|\begin{matrix}0\\0\\0\end{matrix}
\]

Do cột 3 không bán chuẩn hóa được, nên hệ có vô số nghiệm như sau:

Đặt $x_3=b,\ \forall b\in\R$.

Ta có:
\[
\begin{cases}2x_1-x_2=0\\ x_2-2x_3=0\end{cases}
\Rightarrow
\begin{cases}x_1=\dfrac{x_2}{2}=b\\ x_2=2b\end{cases}
\]

\[
\Rightarrow E_{\lambda_2}=E_1=\{X=(x_1,x_2,x_3)=(b,2b,b)\mid b\in\R\}=\{X=b(1,2,1)\mid b\in\R\}
\]

Đặt $u_2=(1,2,1)\neq(0,0,0)=O$

$u_2=(1,2,1)$ là hệ sinh của $E_1$

Ta có: $u_2=(1,2,1)$ là một véc tơ riêng; và $a_2=S_2=\{u_2\}$ là một cơ sở của $E_1$,

và số chiều là $\dim_\R E_{\lambda_2}=1=r_2$ \hfill (2)

\medskip
Với $\lambda_3=4$:

$E_{\lambda_3}=E_4=\{X\in\R^3 \mid (A-\lambda_3I_3)X=O\}$

Ta có dạng ma trận hóa của hệ pt tuyến tính: $(A-\lambda_3I_3)X=O$

\[
\Leftrightarrow
\left[\begin{pmatrix}3 & -1 & 0\\ -1 & 2 & -1\\ 0 & -1 & 3\end{pmatrix}-4\begin{pmatrix}1&0&0\\0&1&0\\0&0&1\end{pmatrix}\ \middle|\ \begin{matrix}0\\0\\0\end{matrix}\right]
\]
\[
\Leftrightarrow
\begin{bmatrix}-1 & -1 & 0\\ -1 & -2 & -1\\ 0 & -1 & -1\end{bmatrix}\big|\begin{matrix}0\\0\\0\end{matrix}
\xrightarrow{(2)\to(2)-(1)}
\begin{bmatrix}-1 & -1 & 0\\ 0 & -1 & -1\\ 0 & -1 & -1\end{bmatrix}\big|\begin{matrix}0\\0\\0\end{matrix}
\xrightarrow{(3)\to(3)-(2)}
\begin{bmatrix}-1 & -1 & 0\\ 0 & -1 & -1\\ 0 & 0 & 0\end{bmatrix}\big|\begin{matrix}0\\0\\0\end{matrix}
\]

Do cột 3 không bán chuẩn hóa được nên hệ có vô số nghiệm như sau:

Đặt $x_3=c,\ \forall c\in\R$

Ta có:
\[
\begin{cases}-x_1-x_2=0\\ -x_2-x_3=0\end{cases}
\Rightarrow
\begin{cases}x_1=-x_2=c\\ x_2=-c\end{cases}
\]

\[
\Rightarrow E_{\lambda_3}=E_4=\{X=(x_1,x_2,x_3)=(c,-c,c)\mid c\in\R\}=\{X=c(1,-1,1)\mid c\in\R\}
\]

Đặt $u_3=(1,-1,1)\neq(0,0,0)=O$

$u_3=(1,-1,1)$ là hệ sinh của $E_4$.

Ta có: $u_3=(1,-1,1)$ là một véc tơ riêng; và $a_3=S_3=\{u_3\}$ là một cơ sở của $E_4$.

và số chiều là $\dim_\R E_{\lambda_3}=1=r_3$ \hfill (3)

\medskip
Từ (1), (2), (3) $\Rightarrow A$ có thể chéo hóa được trên $\R$.

Đặt $a=a_1\cup a_2\cup a_3=\{u_1,u_2,u_3\}$ thì $a$ là một cơ sở của $\R^3$.

Gọi $\beta_0=\{e_1=(1,0,0),\ e_2=(0,1,0),\ e_3=(0,0,1)\}$ là một cơ sở chính tắc của $\R^3$.

Đặt $P=T=T_{\beta_0\to a}=\big([u_1]_{\beta_0}\ \ [u_2]_{\beta_0}\ \ [u_3]_{\beta_0}\big)=\begin{pmatrix}-1 & 1 & 1\\ 0 & 2 & -1\\ 1 & 1 & 1\end{pmatrix}$ thì $T$ là ma trận khả nghịch do $\det(T)=-6\neq 0$ và $A=T\begin{pmatrix}3&0&0\\0&1&0\\0&0&4\end{pmatrix}T^{-1}$.

Ta có: $D=T^{-1}AT=\begin{pmatrix}3&0&0\\0&1&0\\0&0&4\end{pmatrix}$. Đây là dạng chéo hóa cần tìm của $A$.

\[
T^{-1}=\begin{pmatrix}-1/2 & 0 & 1/2\\ 1/6 & 1/3 & 1/6\\ 1/3 & -1/3 & 1/3\end{pmatrix}
=\frac{1}{6}\begin{pmatrix}-3 & 0 & 3\\ 1 & 2 & 1\\ 2 & -2 & 2\end{pmatrix}
\]

\[
\Rightarrow A^m=T\begin{pmatrix}3&0&0\\0&1&0\\0&0&4\end{pmatrix}^{m}T^{-1}
=T\begin{pmatrix}3^m&0&0\\0&1^m&0\\0&0&4^m\end{pmatrix}T^{-1}
\]
\[
=\frac{1}{6}\begin{pmatrix}-1&1&1\\0&2&-1\\1&1&1\end{pmatrix}
\begin{pmatrix}3^m&0&0\\0&1^m&0\\0&0&4^m\end{pmatrix}
\begin{pmatrix}-3&0&3\\1&2&1\\2&-2&2\end{pmatrix}
\]
\[
=\frac{1}{6}\begin{pmatrix}-1\cdot 3^m & 1 & 4^m\\ 0 & 2 & -1\cdot4^m\\ 3^m & 1 & 4^m\end{pmatrix}
\begin{pmatrix}-3&0&3\\1&2&1\\2&-2&2\end{pmatrix}
\]
\[
=\frac{1}{6}\begin{pmatrix}
3\cdot3^m+1+2\cdot4^m & 2-2\cdot4^m & -3\cdot3^m+1+2\cdot4^m\\
2-2\cdot4^m & 4+2\cdot4^m & 2-2\cdot4^m\\
-3\cdot3^m+1+2\cdot4^m & 2-2\cdot4^m & 3\cdot3^m+1+2\cdot4^m
\end{pmatrix},\quad \forall m\in\N.
\]

\medskip
\noindent\textit{\underline{Bài tập tương tự}}: Yêu cầu như ví dụ mẫu

a/ $A=\begin{pmatrix}-4 & -2\\ -4 & 3\end{pmatrix}$. Hãy chéo hóa $A$ và tìm $A^m,\ \forall m\in\N$.

b/ $A=\begin{pmatrix}-5 & -3\\ -3 & 3\end{pmatrix}$. Hãy chéo hóa $A$ và tìm $A^m,\ \forall m\in\N$.

c/ $A=\begin{pmatrix}1 & 0 & 1\\ 0 & 1 & 0\\ 1 & 0 & 1\end{pmatrix}$. Hãy chéo hóa $A$ và tìm $A^m,\ \forall m\in\N$.

d/ $A=\begin{pmatrix}7 & -6 & -10\\ -12 & 17 & 24\\ 12 & -15 & -22\end{pmatrix}$. Hãy chéo hóa $A$ và tìm $A^m,\ \forall m\in\N$.

e/ $A=\begin{pmatrix}8 & -1 & -5\\ -2 & 3 & 1\\ 4 & -1 & -1\end{pmatrix}$. Hãy chéo hóa $A$ và tìm $A^m,\ \forall m\in\N$.

f/ $A=\begin{pmatrix}5 & 6 & -3\\ -1 & 0 & 1\\ 1 & 2 & 1\end{pmatrix}$. Hãy chéo hóa $A$ và tìm $A^m,\ \forall m\in\N$.

g/ $A=\begin{pmatrix}-1 & 3 & -1\\ -3 & 5 & -1\\ -3 & 3 & 1\end{pmatrix}$. Hãy chéo hóa $A$ và tìm $A^m,\ \forall m\in\N$.

h/ $A=\begin{pmatrix}7 & -12 & -2\\ 3 & -4 & 0\\ -2 & 0 & -2\end{pmatrix}$. Hãy chéo hóa $A$ và tìm $A^m,\ \forall m\in\N$.

i/ $A=\begin{pmatrix}2 & -1 & -1\\ -1 & 2 & -1\\ -1 & -1 & 2\end{pmatrix}$. Hãy chéo hóa $A$ và tìm $A^m,\ \forall m\in\N$.

j/ $A=\begin{pmatrix}3 & 2 & 0\\ 2 & 4 & -2\\ 0 & -2 & 5\end{pmatrix}$. Hãy chéo hóa $A$ và tìm $A^m,\ \forall m\in\N$.

\end{document}